% META: {"id": "1NSI-TC-CO-042", "chapitre": "1NSI-TYPES-CONSTRUITS", "type_objet": "corrige", "capacites": ["1NSI-TYPES-CONSTRUITS-C4"], "mode_creation": "ex_nihilo", "sources_inspiration": [], "status": "needs_review"}
% BEGIN-VERIFY
% joueur = {"nom": "Bob", "elo": 1800}
% assert "rang" not in joueur
% assert joueur.get("rang", "inconnu") == "inconnu"
% END-VERIFY
\begin{corrige}{1NSI-TC-EX-042}
\begin{enumerate}
  \item Le programme produit une erreur \lstinline{KeyError: 'rang'}. La clé \lstinline{"rang"} n'existe pas dans le dictionnaire \lstinline{joueur} (qui ne contient que \lstinline{"nom"} et \lstinline{"elo"}). Accéder à une clé inexistante avec la notation \lstinline{dico[cle]} provoque cette erreur.

  \item Deux façons d'éviter l'erreur :

  \textbf{Méthode 1 --- tester l'existence avec \lstinline{in} :}
\begin{python}
if "rang" in joueur:
    print(joueur["rang"])
\end{python}

  \textbf{Méthode 2 --- utiliser \lstinline{get} :}
\begin{python}
print(joueur.get("rang"))
\end{python}
  La méthode \lstinline{get} renvoie \lstinline{None} si la clé n'existe pas, sans provoquer d'erreur.

  \item Pour afficher \lstinline{"inconnu"} par défaut :
\begin{python}
print(joueur.get("rang", "inconnu"))
\end{python}
  Le second argument de \lstinline{get} est la valeur renvoyée si la clé est absente.
\end{enumerate}
\end{corrige}
